\documentclass[11pt]{article}
% The preamble every project's paper shares: % The preamble every project's paper shares: % The preamble every project's paper shares: \input{../../preamble} from
% projects/<project>/paper/. Boards come from chessboard and xskak, plots from
% pgfplots; every number a paper prints comes from its generated.tex.
\usepackage[LSBC4,T1]{fontenc}
\usepackage[utf8]{inputenc}
\usepackage{lmodern}
\usepackage[margin=1in]{geometry}
\usepackage{microtype}
\usepackage{amsmath}
\usepackage{booktabs}
\usepackage{tabularx}
\usepackage{array}
\usepackage{xcolor}
\usepackage[obeyspaces,spaces]{url}
\usepackage{xskak}
\usepackage{pgfplots}
\pgfplotsset{compat=1.18}
\usepackage{caption}
\usepackage{subcaption}
\usepackage{enumitem}
\usepackage[section]{placeins}
\usepackage[hidelinks]{hyperref}

\setchessboard{boardfontencoding=LSBC4, boardfontfamily=skaknew, blackfieldcolor=black!22,
  whitefieldcolor=white, setfontcolors, showmover=false, label=true, labelfontsize=6pt,
  boardfontsize=14pt}

% A move or line in UCI, and a position in FEN, both breakable at a slash.
\newcommand{\uci}[1]{\texttt{#1}}
\newcommand{\fen}[1]{{\footnotesize\path{#1}}}
\newcolumntype{L}{>{\raggedright\arraybackslash}X}

\setlength{\parskip}{0.4em}
\setlength{\parindent}{0pt}
 from
% projects/<project>/paper/. Boards come from chessboard and xskak, plots from
% pgfplots; every number a paper prints comes from its generated.tex.
\usepackage[LSBC4,T1]{fontenc}
\usepackage[utf8]{inputenc}
\usepackage{lmodern}
\usepackage[margin=1in]{geometry}
\usepackage{microtype}
\usepackage{amsmath}
\usepackage{booktabs}
\usepackage{tabularx}
\usepackage{array}
\usepackage{xcolor}
\usepackage[obeyspaces,spaces]{url}
\usepackage{xskak}
\usepackage{pgfplots}
\pgfplotsset{compat=1.18}
\usepackage{caption}
\usepackage{subcaption}
\usepackage{enumitem}
\usepackage[section]{placeins}
\usepackage[hidelinks]{hyperref}

\setchessboard{boardfontencoding=LSBC4, boardfontfamily=skaknew, blackfieldcolor=black!22,
  whitefieldcolor=white, setfontcolors, showmover=false, label=true, labelfontsize=6pt,
  boardfontsize=14pt}

% A move or line in UCI, and a position in FEN, both breakable at a slash.
\newcommand{\uci}[1]{\texttt{#1}}
\newcommand{\fen}[1]{{\footnotesize\path{#1}}}
\newcolumntype{L}{>{\raggedright\arraybackslash}X}

\setlength{\parskip}{0.4em}
\setlength{\parindent}{0pt}
 from
% projects/<project>/paper/. Boards come from chessboard and xskak, plots from
% pgfplots; every number a paper prints comes from its generated.tex.
\usepackage[LSBC4,T1]{fontenc}
\usepackage[utf8]{inputenc}
\usepackage{lmodern}
\usepackage[margin=1in]{geometry}
\usepackage{microtype}
\usepackage{amsmath}
\usepackage{booktabs}
\usepackage{tabularx}
\usepackage{array}
\usepackage{xcolor}
\usepackage[obeyspaces,spaces]{url}
\usepackage{xskak}
\usepackage{pgfplots}
\pgfplotsset{compat=1.18}
\usepackage{caption}
\usepackage{subcaption}
\usepackage{enumitem}
\usepackage[section]{placeins}
\usepackage[hidelinks]{hyperref}

\setchessboard{boardfontencoding=LSBC4, boardfontfamily=skaknew, blackfieldcolor=black!22,
  whitefieldcolor=white, setfontcolors, showmover=false, label=true, labelfontsize=6pt,
  boardfontsize=14pt}

% A move or line in UCI, and a position in FEN, both breakable at a slash.
\newcommand{\uci}[1]{\texttt{#1}}
\newcommand{\fen}[1]{{\footnotesize\path{#1}}}
\newcolumntype{L}{>{\raggedright\arraybackslash}X}

\setlength{\parskip}{0.4em}
\setlength{\parindent}{0pt}

% Written by projects/period-writer/paper/generate.mjs; do not edit by hand.

\newcommand{\AppCommit}{fa07ea1ed8}
\newcommand{\Cases}{nine}
\newcommand{\Weeks}{seven}
\newcommand{\Months}{two}
\newcommand{\TrainCases}{week 35, week 37, week 39, week 41 and August}
\newcommand{\HeldCases}{week 36, week 38, week 40 and September}
\newcommand{\TrainN}{five}
\newcommand{\HeldN}{four}
\newcommand{\WeekGamesMin}{19}
\newcommand{\WeekGamesMax}{59}
\newcommand{\MonthGames}{307 and 172}
\newcommand{\FrozenBriefs}{four}
\newcommand{\CaseSpan}{weeks 35 to 41 and the months of August and September 2026}
\newcommand{\CasesCap}{Nine}
\newcommand{\WriterRuns}{25}
\newcommand{\FirstRuns}{18}
\newcommand{\SecondPairCases}{week 35, August and week 36}
\newcommand{\SecondPairN}{three}
\newcommand{\BaseOnlySecond}{week 37}
\newcommand{\MaxReps}{two}
\newcommand{\CandPairs}{twelve}
\newcommand{\NullPairs}{four}
\newcommand{\NullCases}{week 35, week 37, August and week 36}
\newcommand{\Spent}{225}
\newcommand{\Cap}{250}
\newcommand{\WriterSpend}{168}
\newcommand{\GraderSpend}{58}
\newcommand{\SevenDayFirst}{12}
\newcommand{\SevenDayLast}{25}
\newcommand{\BaseWeekCost}{6.62}
\newcommand{\BaseMonthCost}{5.78}
\newcommand{\BaseMinutes}{21}
\newcommand{\BaseRounds}{5.5}
\newcommand{\BasePassed}{13}
\newcommand{\BaseRuns}{13}
\newcommand{\BaseBoards}{44}
\newcommand{\BaseMonthBoards}{0, 0 and 1}
\newcommand{\CandWeekCost}{7.55}
\newcommand{\CandMonthCost}{5.49}
\newcommand{\CandMinutes}{20}
\newcommand{\CandRounds}{6.3}
\newcommand{\CandPassed}{12}
\newcommand{\CandRuns}{12}
\newcommand{\CandBoards}{42}
\newcommand{\CandMonthBoards}{1, 1 and 2}
\newcommand{\BaseMonthRuns}{three}
\newcommand{\BaseExhibits}{44}
\newcommand{\BaseFullOrder}{32}
\newcommand{\BaseThreat}{68}
\newcommand{\BaseRefutation}{70}
\newcommand{\BaseKeeps}{73}
\newcommand{\CandExhibits}{42}
\newcommand{\CandFullOrder}{79}
\newcommand{\CandThreat}{88}
\newcommand{\CandRefutation}{79}
\newcommand{\CandKeeps}{83}
\newcommand{\TrainTruth}{+0.20 {\footnotesize ($-$0.40 to +0.80)}}
\newcommand{\TrainExplanation}{+0.40 {\footnotesize ($-$0.40 to +1.00)}}
\newcommand{\TrainDuties}{+0.10 {\footnotesize ($-$0.50 to +0.60)}}
\newcommand{\TrainPlain}{0.00 {\footnotesize (0.00 to 0.00)}}
\newcommand{\TrainOverall}{+0.30 {\footnotesize ($-$0.40 to +0.90)}}
\newcommand{\TrainNullExplanation}{$-$0.33 {\footnotesize ($-$1.00 to 0.00)}}
\newcommand{\HeldTruth}{0.00 {\footnotesize (0.00 to 0.00)}}
\newcommand{\HeldExplanation}{+0.38 {\footnotesize (0.00 to +0.75)}}
\newcommand{\HeldDuties}{$-$0.13 {\footnotesize ($-$0.38 to 0.00)}}
\newcommand{\HeldPlain}{0.00 {\footnotesize (0.00 to 0.00)}}
\newcommand{\HeldOverall}{+0.13 {\footnotesize (0.00 to +0.38)}}
\newcommand{\HeldNullExplanation}{+1.00 {\footnotesize (+1.00 to +1.00)}}
\newcommand{\AllTruth}{+0.11 {\footnotesize ($-$0.22 to +0.44)}}
\newcommand{\AllExplanation}{+0.39 {\footnotesize ($-$0.06 to +0.78)}}
\newcommand{\AllDuties}{0.00 {\footnotesize ($-$0.33 to +0.33)}}
\newcommand{\AllPlain}{0.00 {\footnotesize (0.00 to 0.00)}}
\newcommand{\AllOverall}{+0.22 {\footnotesize ($-$0.17 to +0.56)}}
\newcommand{\AllNullExplanation}{0.00 {\footnotesize ($-$0.75 to +0.75)}}
\newcommand{\CandDecided}{eight}
\newcommand{\ExplWins}{seven}
\newcommand{\ExplLosses}{one}
\newcommand{\ExplTies}{four}
\newcommand{\OverallWins}{five}
\newcommand{\OverallLosses}{one}
\newcommand{\OverallTies}{six}
\newcommand{\Decided}{130}
\newcommand{\SecondShare}{63}
\newcommand{\NullDecisive}{two}
\newcommand{\SameN}{nine}
\newcommand{\SameBroken}{zero}
\newcommand{\DecoyN}{nine}
\newcommand{\DecoyExpl}{zero}
\newcommand{\DecoyLost}{four}
\newcommand{\DecoyRatio}{1.21}
\newcommand{\RuleTable}{%
\begin{tabular}{lrl}
\toprule
Condition & Estimate (95\% interval) & Holds \\
\midrule
Explanation on train above 0 & +0.40 {\footnotesize ($-$0.40 to +1.00)} & \textbf{no} \\
Explanation held out above 0 & +0.38 {\footnotesize (0.00 to +0.75)} & \textbf{no} \\
Overall not below 0 & +0.22 {\footnotesize ($-$0.17 to +0.56)} & yes \\
False claims not up & +0.78 {\footnotesize ($-$0.56 to +2.22)} & yes \\
False board claims not up & +1.22 {\footnotesize ($-$1.11 to +3.56)} & yes \\
Check passes as often & -- & yes \\
Controls clean & -- & yes \\
\bottomrule
\end{tabular}
}
\newcommand{\RuleFailed}{explanation on train above 0 and explanation held out above 0}
\newcommand{\RuleFailedN}{two}
\newcommand{\HeldTieCases}{week 40 and September}
\newcommand{\HeldTieN}{two}
\newcommand{\TrainLossCases}{week 37}
\newcommand{\TrainFalseDiff}{$-$0.8 {\footnotesize ($-$1.6 to +0.2)}}
\newcommand{\HeldFalseDiff}{+2.8 {\footnotesize (+1.5 to +4.0)}}
\newcommand{\BaseFalse}{2.1}
\newcommand{\BaseUnsupported}{9.9}
\newcommand{\BaseMisleading}{5.5}
\newcommand{\BaseBoardFalse}{2.5}
\newcommand{\CandFalse}{2.9}
\newcommand{\CandUnsupported}{8.5}
\newcommand{\CandMisleading}{4.8}
\newcommand{\CandBoardFalse}{3.6}
\newcommand{\FalseDiff}{+0.8 {\footnotesize ($-$0.6 to +2.2)}}
\newcommand{\BoardDiff}{+1.2 {\footnotesize ($-$1.1 to +3.6)}}
\newcommand{\KindsTable}{%
\begin{tabular}{lrrrr}
\toprule
 & \multicolumn{2}{c}{Today's prompt} & \multicolumn{2}{c}{Candidate} \\
\cmidrule(lr){2-3}\cmidrule(lr){4-5}
 & & Unsupported & & Unsupported \\
What the claim is about & False & or misleading & False & or misleading \\
\midrule
Who stood better, the best move & 0.5 & 1.1 & 0.9 & 1.3 \\
A skill, metric or focus reading & 0.3 & 3.1 & 0.6 & 1.9 \\
Habits and trends & 0.2 & 2.3 & 0.5 & 2.1 \\
Results, records, ratings & 0.4 & 1.5 & 0.3 & 1.2 \\
What an earlier review said & 0.2 & 2.1 & 0.3 & 1.6 \\
What a piece does on the board & 0.2 & 1.0 & 0.3 & 0.7 \\
Counts and tallies & 0.2 & 1.8 & 0.0 & 1.0 \\
Other & 0.0 & 0.5 & 0.1 & 0.4 \\
The clock & 0.1 & 0.6 & 0.0 & 0.4 \\
\midrule
Every kind & 2.0 & 13.9 & 2.9 & 10.5 \\
\bottomrule
\end{tabular}
}
\newcommand{\TopKind}{who stood better, the best move}
\newcommand{\SecondKind}{a skill, metric or focus reading}
\newcommand{\ThirdKind}{results, records, ratings}
\newcommand{\TopThreeShare}{61}
\newcommand{\HeldShare}{13}
\newcommand{\HeldBoard}{33}
\newcommand{\WalkReviewsCap}{Seven}
\newcommand{\WalkReviews}{seven}
\newcommand{\WalkOf}{15}
\newcommand{\WalkFirst}{36}
\newcommand{\WalkLast}{41}
\newcommand{\WalkBase}{five}
\newcommand{\WalkBaseOf}{eight}
\newcommand{\WalkCand}{two}
\newcommand{\WalkCandOf}{seven}
\newcommand{\WalkFlagged}{five}
\newcommand{\BaseFalseNotHeld}{2.0}
\newcommand{\BaseLineFalse}{0.15}
\newcommand{\CandFalseNotHeld}{2.9}
\newcommand{\CandLineFalse}{1.00}
\newcommand{\ExtraFalse}{0.9}
\newcommand{\ExtraLines}{0.8}
\newcommand{\BasePlayed}{40}
\newcommand{\BasePlayedWrong}{seven}
\newcommand{\BasePlayedReviews}{six}
\newcommand{\BasePlayedShare}{18}
\newcommand{\CandPlayed}{36}
\newcommand{\CandPlayedWrong}{eleven}
\newcommand{\CandPlayedReviews}{seven}
\newcommand{\CandPlayedShare}{31}
\newcommand{\LandedPlayed}{27}
\newcommand{\LandedPlayedWrong}{two}
\newcommand{\LandedPlayedReviews}{two}
\newcommand{\LandedPlayedShare}{7}
\newcommand{\PlayedAll}{76}
\newcommand{\PlayedWrongAll}{18}
\newcommand{\PlayedReviewsAll}{13}
\newcommand{\PlayedShareAll}{24}
\newcommand{\LandedReviews}{seven}
\newcommand{\PlayedCheckPR}{125}
\newcommand{\CaseTable}{%
\begin{tabular}{llrrrcc}
\toprule
 & & & \multicolumn{2}{c}{Exhibits in order} & \multicolumn{2}{c}{Judge, per pair} \\
Period & Split & Games & Today & Candidate & Explanation & Overall \\
\midrule
Week 35 & train & 55 & 3/9 & 7/9 & + + & + + \\
Week 37 & train & 19 & 4/9 & 3/4 & $-$ & $-$ \\
Week 39 & train & 59 & 1/4 & 3/4 & + & + \\
Week 41 & train & 34 & 3/4 & 3/4 & 0 & 0 \\
August & train & 307 & 0/0 & 2/2 & + + & 0 + \\
Week 36 & held out & 32 & 1/8 & 6/8 & 0 + & 0 + \\
Week 38 & held out & 49 & 0/4 & 3/4 & + & 0 \\
Week 40 & held out & 33 & 2/5 & 4/5 & 0 & 0 \\
September & held out & 172 & 0/1 & 2/2 & 0 & 0 \\
\bottomrule
\end{tabular}
}


\title{Does the books' order make the weekly review explain better?\\[0.3em]
\large One sentence added to the weekly and monthly routines' prompts, replayed on \Cases{} past
periods}
\author{Chess Coach research lab}
\date{}

\begin{document}
\maketitle

\begin{abstract}
The app's weekly and monthly routines have a model write a review from a brief the app's services
compute. Two chess books explain a position in one order: the threat, then why the tempting move
fails, then what the better move keeps. One sentence asking for that order was added to the live
prompts, and both prompts rewrote \Cases{} past periods (\Weeks{} weeks and \Months{} months)
offline with the app's code at commit \texttt{\AppCommit}: \WriterRuns{} reviews, for \$\Spent{} of
model calls. The sentence changed the writing: \CandFullOrder\% of the candidate's board exhibits
gave all three parts in order, against \BaseFullOrder\% under today's prompt. A blind judge, reading
each pair in both orders, preferred the candidate's explanations in \ExplWins{} of \CandPairs{} pairs
and today's in \ExplLosses, but its interval reached zero on the train periods and on the held-out
ones, so the rule fixed before the runs keeps today's prompts. The same judge sided with whichever
review it read second in \SecondShare\% of the verdicts that picked a side, and two runs of today's
prompt on one period split on explanation in \NullDecisive{} of \NullPairs{} periods. On the held-out
periods the candidate's reviews made more false claims, and most of its extra false claims quote a
line of moves. Two errors do not depend on the prompt. Of the \PlayedAll{} lines the reviews tag
\emph{Played}, \PlayedWrongAll{} are not what the game played, and the app's checker, which replays
each line only for legality, lets them through. And the reviews call a rule about the king in a won
endgame ``the pawn break'', because the rule ledger files it under that concept.
\end{abstract}

\section{The question}

Levy Rozman's \emph{How to Win at Chess} (2023) and \emph{How to Win at Chess: Next Level} (2026)
answered twenty questions about the app's coaching. Asked how a coach should explain a better move,
they gave an order their worked examples share: first the threat on the board, then why the move
played (or the tempting one) fails, with the reply that refutes it, then what the better move keeps.
They do not prescribe one sentence order or a length, and their examples start in different places,
so the candidate tests the order they share most, as one instruction added to step 4 of the weekly
and monthly prompts, right after the sentence that says to write the annotations by filling the
brief's slots:

\begin{quote}\small
Each board that shows a better move explains it in its prose in this order: first the threat on the
board, the piece and the square it aims at; then why the move played, or the tempting one, fails,
with the reply that refutes it; then what the better move keeps, the square, piece or line it holds.
\end{quote}

Nothing else in the prompts changes. Measured beside the comparison, for both prompts: which false
claims the weekly and monthly reviews make.

\section{The cases and the writer}

The cases are the \Cases{} periods with a review on the app's main branch from its command-line
era: \CaseSpan. Each is a checkout of the app at
\texttt{\AppCommit}, with the coach data as main held it just before that period's review landed,
and the period's brief prepared again by the app's own \texttt{coach weekly prepare} or
\texttt{coach monthly prepare}, offline. A week holds \WeekGamesMin{} to \WeekGamesMax{} games and
the months \MonthGames. The split was fixed before any run, interleaved by date: train is
\TrainCases; held out is \HeldCases.

The writer runs as the routines do: Opus 5.5, headless, with nothing inherited from the session
that launched it, and a short preamble that names the period, says the brief is ready and stops the
writer once the strict check passes (no video, publish or commit). Each prompt wrote every case once,
and again where the budget allowed a second pair (\SecondPairCases); today's prompt also wrote
\BaseOnlySecond{} twice. Run $k$ of the candidate is judged against run $k$ of today's prompt, and today's two runs
of a period against each other give the judge's noise floor (\NullCases).

\section{The graders}

A script with no model reads each run: whether the final strict check passes, every check the
writer ran, commands it should not have run, cost and minutes. Sonnet 5.5 audits each review alone.
It reads the evidence (the brief without its bulk, a line per game, the clock's turning points, the
earlier reviews the period's review audits, and for every board exhibit the position as the archive
replays it, what was played there and next, the engine's best move and reply, and what each piece
attacks and guards before and after each move the exhibit names) and lists the claims the evidence
contradicts (\emph{false}), does not carry (\emph{unsupported}) or carries but slants
(\emph{misleading}). For each exhibit it says whether the prose names the threat, the refuting reply
and what the better move keeps, in that order, and lists its false claims about the board. Its
prompts were settled on one run of today's prompt before any candidate run was graded. Opus 5.5,
tried as the auditor on that run, did not list the same false claims, so a review's count is a floor
on what it gets wrong, compared like for like.

Sonnet 5.5 also judges each pair blind, in both orders, on five criteria: truth, explanation, the
prompt's duties, plain words, and overall. A side wins a criterion only when both orders agree.
Two controls check the judge on every case: a review of today's prompt against itself, and against
a decoy whose exhibits restate what they already say at the candidate's length, with no square,
move, piece or number the original does not name.

\section{The rule}

A pair scores $+1$ when the candidate wins a criterion, $-1$ when it loses and $0$ for a tie, and a
case scores the mean over its pairs; intervals are 95\% percentile bootstrap intervals over cases.
The candidate replaces the live prompts only when every condition holds
(Table~\ref{tab:rule}): its explanation interval lies above zero on the train cases and on the
held-out ones, the overall interval does not lie below zero, neither the false claims nor the false
board claims per review rise by an interval above zero, the strict check passes as often, and the
controls are clean. With \HeldN{} or \TrainN{} cases a split, the first two need the candidate to
win nearly every case; that is the bar. The \RuleFailedN{} that fail are \RuleFailed, so today's prompts stay.

\begin{table}[h]
\centering
\small\RuleTable
\caption{The rule, fixed before the candidate's first run. Explanation and overall are the judge's
scores per case; the claim rows are the audit's counts per review, candidate minus today's prompt,
over all \Cases{} cases.}
\label{tab:rule}
\end{table}

\section{What the sentence changed}

The writer did what the sentence asks. Pooled over the exhibits, \CandFullOrder\% of the
candidate's \CandExhibits{} gave all three parts in order, against \BaseFullOrder\% of today's
\BaseExhibits. The threat was named in \CandThreat\% (today \BaseThreat\%), the refuting reply in
\CandRefutation\% (\BaseRefutation\%), and what the better move keeps in \CandKeeps\%
(\BaseKeeps\%). Every review passed the strict check in the end, \BasePassed{} of \BaseRuns{} and
\CandPassed{} of \CandRuns, after \BaseRounds{} and \CandRounds{} check rounds on average. A weekly
review cost \$\BaseWeekCost{} under today's prompt and \$\CandWeekCost{} under the candidate, and a
monthly one \$\BaseMonthCost{} and \$\CandMonthCost; a review took about \BaseMinutes{} minutes
either way. The monthly reviews show few boards: today's \BaseMonthRuns{} carried
\BaseMonthBoards{} and the candidate's \CandMonthBoards, so the explanation criterion rests on the weeks
(Table~\ref{tab:cases}).

\begin{table}[h]
\centering
\small\CaseTable
\caption{Each period: its games, the exhibits that give all three parts in order (over all of a
prompt's runs of the period), and the judge's verdict on each pair, from the candidate's side ($+$ it
wins, $-$ it loses, $0$ a tie).}
\label{tab:cases}
\end{table}

\section{What the judge saw}

On explanation the judge preferred the candidate in \ExplWins{} of the \CandPairs{} pairs, today's
prompt in \ExplLosses{} and neither in \ExplTies. That is \TrainExplanation{} over the train cases
and \HeldExplanation{} over the held-out ones. Today's prompt took \TrainLossCases{} and
\HeldTieCases{} tied, so neither interval clears zero. Overall the candidate won \OverallWins{}
pairs, lost \OverallLosses{} and tied \OverallTies{} (\AllOverall{} over all \Cases{} cases), and
the intervals of truth (\AllTruth), duties (\AllDuties) and plain words (\AllPlain) hold zero.

Two things say the judge is a weak instrument for this question. Today's prompt against itself,
two runs of one period, gave a decisive explanation verdict in \NullDecisive{} of \NullPairs{}
periods (\AllNullExplanation), against \CandDecided{} of \CandPairs{} for the candidate's pairs.
And across all pairs
and criteria, \SecondShare\% of the \Decided{} single-order verdicts that picked a side picked the
review it read second; requiring both orders to agree turns most of those into ties. The controls
are clean: the same review on both sides tied on every criterion in \SameN{} of \SameN{} cases, and
the decoy, whose exhibits ran \DecoyRatio{} times as long as the original's, won explanation in
\DecoyExpl{} of \DecoyN{} cases and lost overall in \DecoyLost. The judge does not reward length
alone, but with \Cases{} cases it cannot tell this sentence from the writer's run-to-run variation.

\section{What the reviews get wrong}

Per review, today's prompt made \BaseFalse{} false claims, \BaseUnsupported{} unsupported and
\BaseMisleading{} misleading ones, and \BaseBoardFalse{} false board claims; the candidate made
\CandFalse, \CandUnsupported, \CandMisleading{} and \CandBoardFalse. Candidate minus today's, the
false claims are \TrainFalseDiff{} on the train cases and \HeldFalseDiff{} on the held-out ones,
where every case went up; over all \Cases{} they are \FalseDiff, inside rule 4.

A second blind Sonnet call sorted every listed claim by what it is about, and marked an entry whose
own evidence says the claim holds: \HeldShare\% of the claims listed in words, and \HeldBoard\% of
the false board claims, which Table~\ref{tab:kinds} leaves out. The false claims are mostly about
\TopKind, \SecondKind{} and \ThirdKind{} (\TopThreeShare\% of them); the unsupported ones spread over
the brief's readings, trends and what earlier reviews said, where the evidence the auditor read may
lack what the writer read.

\begin{table}[h]
\centering
\small\KindsTable
\caption{Claims per review the audit lists, by what they are about, without those whose evidence
says they hold.}
\label{tab:kinds}
\end{table}

\textbf{Lines of moves.} False claims that quote a line of two numbered moves or more run
\BaseLineFalse{} a review under today's prompt and \CandLineFalse{} under the candidate: \ExtraLines{}
of the candidate's \ExtraFalse{} extra false claims a review. The sentence asks for the reply that
refutes the move; the writer answered with lines, and wrote some of them wrong.

\textbf{Lines tagged Played.} A board exhibit carries lines with a tag, most often \emph{Played} and
\emph{Better}. A script with no model replayed every Played line from the exhibit's position and
compared it with the game in the archive. Of the \PlayedAll{} Played lines, \PlayedWrongAll{} leave
the game (\BasePlayedShare\% under today's prompt, \CandPlayedShare\% under the candidate), in
\PlayedReviewsAll{} of the \WriterRuns{} reviews, and so do \LandedPlayedWrong{} of the
\LandedPlayed{} in the reviews that landed. The usual shape is the move played followed by the
engine's refutation, or a later move the game never had. The strict check replays each line for
legality from the board's position but does not compare a Played line with the game, so none of
these was refused.

\textbf{The king walk called the pawn break.} The weekly review reports the rule ledger's budget,
which caps the live rules per concept. In weeks \WalkFirst{} to \WalkLast, the only live rule under
the concept ``Pawn structure and the break'' is about the king in a won endgame. Of the \WalkOf{}
weekly reviews of those weeks, \WalkReviews{} (\WalkBase{} of \WalkBaseOf{} under
today's prompt, \WalkCand{} of \WalkCandOf{} under the candidate) count it as ``the pawn break'', and
the auditor marks \WalkFlagged{} of them false; the rest say ``pawn structure''. The writer follows
the prompt, which names a concept by its name. The claim is false because the ledger files the rule
under a concept it is not about.

\section{What the app takes}

The live weekly and monthly prompts stay. The app took both fixes that do not depend on the
prompt: its strict check now refuses a line tagged Played that leaves the game from the board's
position, for the weeks and months after the ones these cases cover (chess PR \PlayedCheckPR), and
its rule ledger files the king-walk rule under converting, so the per-concept count no longer puts
it under the pawn break (chess PR \WalkRefilePR). If the order is tried again, it would want the
engine's refutation in the brief for each exhibit, so the line the writer shows is one the engine
played out (an inference from the lines above, not tested), and a judge that compares the same
exhibit under both prompts, which needs the exhibits fixed before the writer runs.

\section{Limits}

\CasesCap{} cases, and at most \MaxReps{} runs of a prompt on one; the judge's verdicts on the same prompt vary
nearly as much as between prompts, and it leans to the review it reads second. The audit's counts
are a floor, and a share of what it lists holds by its own evidence. One writer model wrote every
review, offline: the level falls back to the profile synced at the time, and the brief is the one
the services compute at \texttt{\AppCommit}, not always the one the period's routine read. The
writer eval spent \$\WriterSpend{} and the graders \$\GraderSpend{} at list price; on the owner's
subscription the seven-day usage window moved from \SevenDayFirst\% to \SevenDayLast\% while it ran.

\end{document}
